\chapter{FPU: What an Arrow Is Worth Before You Have Ever Taken It}
\label{ch:fpu_transition}

This is the part you pushed hardest on, and you were right.

\begin{studentq}
``Measured quality $Q$ is for the state prior to the move Kd6, it is not the state that Kd6 caused\dots\ The bonus is added to the prior value of $s$ and not the current state after the move Kd6.''
\end{studentq}

Correct, and here is the precise statement.

\section{Phase 1: Arrow Never Taken ($n = 0$)}

There is no child position. The network has evaluated nothing down there. There is literally no measurement to average. So the engine needs a placeholder, and it uses the \textbf{parent position's own value} as the starting floor:
\begin{equation}
  Q(a) = \underbrace{Q(s)}_{\substack{\text{parent's} \\ \text{running average}}} - \underbrace{k \cdot \sqrt{\sum_{\text{visited}} P(a')}}_{\substack{\text{the FPU reduction} \\ (k = 0.33)}}
\end{equation}

Two details that are easy to get wrong, and both matter:
\begin{itemize}[leftmargin=2em]
  \item The first term is the parent's \textbf{running average $Q(s)$}, not its original network value $V(s)$. Early on they are equal, because the parent's only sample is its own evaluation. They drift apart as children report back.
  \item The sum is over the \textbf{policy mass already explored}, not the number of moves explored. Spending a visit on a 45\% move makes the engine far more cautious about the rest than spending one on a 5\% move would.
\end{itemize}

Read the reduction physically: \textit{``we have already spent effort on arrows covering this much of the network's attention, and none of it beat the parent, so be a little more pessimistic about the untried ones.''} At the very start nothing has been visited, the sum is 0, and the placeholder is simply the parent's value.

This is called \textbf{First Play Urgency} (FPU).

\section{Phase 2: The Moment the Arrow Is Taken ($n = 0 \to 1$)}

\begin{studentq}
``When Kd6 is visited once and its state evaluation is made, does its previous default value get replaced by the new value?'' --- \textbf{Yes. Completely.}
\end{studentq}

\begin{align*}
  \text{\textbf{BEFORE:}} \quad & Q(\text{Kd6}) = 0.97602 \quad \text{(borrowed from parent; a temporary placeholder)} \\[1ex]
  & \text{The search plays Kd6, evaluates child, returns } x_1 = +0.96766 \\[1ex]
  \text{\textbf{AFTER:}} \quad & Q(\text{Kd6}) = \mathbf{0.96766} \quad \text{(its own measurement; } n = 1, W = 0.96766\text{)}
\end{align*}

\begin{center}
\begin{minipage}{\linewidth}
\centering
% fig_c_fpu.tex -- The FPU placeholder to measured reality transition
\begin{tikzpicture}[scale=0.90, every node/.style={font=\small}]
  % Panel A: Before First Visit (n = 0)
  \begin{scope}[xshift=0cm]
    \node[font=\footnotesize\bfseries, text=NavyBlue] at (1.8, 3.7) {%
      Panel A: Before First Visit ($n = 0$)%
    };
    
    % Dashed box representing borrowed / hollow placeholder
    \draw[NavyBlue, line width=0.9pt, dashed] (0.6, 0.6) rectangle (3.0, 3.0);
    \node[font=\tiny\bfseries, text=NavyBlue, align=center] at (1.8, 2.1) {%
      $Q_{\text{FPU}} = +0.97602$\\[0.4ex]
      \scriptsize\textcolor{WorkGrey}{(Dashed / Hollow)}%
    };
    \node[font=\tiny\bfseries, text=WorkGrey] at (1.8, 1.2) {Borrowed from Parent};
    
    \node[align=center, font=\scriptsize, text=WorkGrey] at (1.8, -0.2) {%
      \textbf{A Placeholder Guess}\\
      No child exists yet.\\
      No measurement has occurred.%
    };
  \end{scope}

  % Transition arrow
  \draw[GoldPath, line width=1.5pt, -{Stealth[length=3.5mm]}] (3.5, 1.8) -- (5.0, 1.8);
  \node[font=\tiny\bfseries, text=GoldPath, align=center] at (4.25, 2.3) {Visit 1\\happens!};

  % Panel B: After First Visit (n = 1)
  \begin{scope}[xshift=5.8cm]
    \node[font=\footnotesize\bfseries, text=NavyBlue] at (3.0, 3.7) {%
      Panel B: After First Visit ($n = 1$)%
    };
    
    % Old placeholder crossed out
    \draw[WorkGrey!40, line width=0.6pt, dashed] (0.6, 0.6) rectangle (2.2, 3.0);
    \draw[PitfallRed, line width=1.2pt] (0.6, 0.6) -- (2.2, 3.0);
    \draw[PitfallRed, line width=1.2pt] (0.6, 3.0) -- (2.2, 0.6);
    \node[font=\tiny\bfseries, text=PitfallRed, fill=white, inner sep=1pt, align=center] at (1.4, 1.8) {0.97602\\DISCARDED!};
    
    % Solid box representing new measured value
    \fill[BuildBlue!70] (2.7, 0.6) rectangle (4.3, 3.0);
    \draw[NavyBlue, line width=1.0pt] (2.7, 0.6) rectangle (4.3, 3.0);
    \node[font=\tiny\bfseries, text=white, align=center] at (3.5, 2.0) {$Q = +0.96766$};
    \node[font=\tiny\bfseries, text=white] at (3.5, 1.4) {Solid Measured};
    
    \node[align=center, font=\scriptsize, text=NavyBlue] at (2.5, -0.2) {%
      \textbf{Permanent Measured Reality}\\
      Leaf returned $x_1 = +0.96766$.\\
      $Q = W/n = 0.96766/1$.\\[0.3ex]
      \textbf{\textcolor{BuildBlue}{Replaced completely, never blended!}}%
    };
  \end{scope}
\end{tikzpicture}

\captionof{figure}{\textbf{The FPU transition.} The placeholder is discarded entirely upon the first visit; it is never blended with real measurements.}
\label{fig:fpu}
\end{minipage}
\end{center}

The placeholder is \textbf{discarded, not blended}. It was never data. From this moment on, $Q(\text{Kd6})$ is determined entirely by real evaluations returning from below that arrow.

And crucially: the three arrows still at $n = 0$ keep using the placeholder --- now with a slightly larger FPU reduction, because \texttt{Kd6}'s prior has joined the ``already visited'' sum.

\begin{keyidea}[title={One Line to Remember}]
\begin{align*}
  n = 0 \quad &\longrightarrow \quad Q \text{ comes from the \textbf{PARENT} position (borrowed, temporary)} \\
  n \ge 1 \quad &\longrightarrow \quad Q \text{ comes from the \textbf{CHILD} positions (measured, permanent)}
\end{align*}
\end{keyidea}
